\documentclass[10pt,a4paper]{amsart}
\usepackage{amsmath,amssymb,amsthm,graphicx}
\usepackage{verbatim}
\newenvironment{reference}{\comment}{\endcomment}
%\newenvironment{reference}{}{}
\newenvironment{slogan}{\comment}{\endcomment}
\newenvironment{history}{\comment}{\endcomment}

% For commutative diagrams we use Xy-pic
\usepackage[all]{xy}

% We use 2cell for 2-commutative diagrams.
\xyoption{2cell}
% Theorem environments.
%
\theoremstyle{plain}
\newtheorem{theorem}[subsection]{Theorem}
\newtheorem{proposition}[subsection]{Proposition}
\newtheorem{lemma}[subsection]{Lemma}

\theoremstyle{definition}
\newtheorem{definition}[subsection]{Definition}
\newtheorem{example}[subsection]{Example}
\newtheorem{exercise}[subsection]{Exercise}
\newtheorem{situation}[subsection]{Situation}

\theoremstyle{remark}
\newtheorem{remark}[subsection]{Remark}
\newtheorem{remarks}[subsection]{Remarks}

\numberwithin{equation}{subsection}

% Macros
%
\def\lim{\mathop{\mathrm{lim}}\nolimits}
\def\colim{\mathop{\mathrm{colim}}\nolimits}
\def\Spec{\mathop{\mathrm{Spec}}}
\def\Hom{\mathop{\mathrm{Hom}}\nolimits}
\def\Ext{\mathop{\mathrm{Ext}}\nolimits}
\def\SheafHom{\mathop{\mathcal{H}\!\mathit{om}}\nolimits}
\def\SheafExt{\mathop{\mathcal{E}\!\mathit{xt}}\nolimits}
\def\Sch{\mathit{Sch}}
\def\Mor{\mathop{\mathrm{Mor}}\nolimits}
\def\Ob{\mathop{\mathrm{Ob}}\nolimits}
\def\Sh{\mathop{\mathit{Sh}}\nolimits}
\def\NL{\mathop{N\!L}\nolimits}
\def\CH{\mathop{\mathrm{CH}}\nolimits}
\def\proetale{{pro\text{-}\acute{e}tale}}
\def\etale{{\acute{e}tale}}
\def\QCoh{\mathit{QCoh}}
\def\Ker{\mathop{\mathrm{Ker}}}
\def\Im{\mathop{\mathrm{Im}}}
\def\Coker{\mathop{\mathrm{Coker}}}
\def\Coim{\mathop{\mathrm{Coim}}}

% Boxtimes
%
\DeclareMathSymbol{\boxtimes}{\mathbin}{AMSa}{"02}

%
% Macros for moduli stacks/spaces
%
\def\QCohstack{\mathcal{QC}\!\mathit{oh}}
\def\Cohstack{\mathcal{C}\!\mathit{oh}}
\def\Spacesstack{\mathcal{S}\!\mathit{paces}}
\def\Quotfunctor{\mathrm{Quot}}
\def\Hilbfunctor{\mathrm{Hilb}}
\def\Curvesstack{\mathcal{C}\!\mathit{urves}}
\def\Polarizedstack{\mathcal{P}\!\mathit{olarized}}
\def\Complexesstack{\mathcal{C}\!\mathit{omplexes}}
% \Pic is the operator that assigns to X its picard group, usage \Pic(X)
% \Picardstack_{X/B} denotes the Picard stack of X over B
% \Picardfunctor_{X/B} denotes the Picard functor of X over B
\def\Pic{\mathop{\mathrm{Pic}}\nolimits}
\def\Picardstack{\mathcal{P}\!\mathit{ic}}
\def\Picardfunctor{\mathrm{Pic}}
\def\Deformationcategory{\mathcal{D}\!\mathit{ef}}
\usepackage{cleveref}
\crefname{originallemma}{original lemma}{original lemmas}
\Crefname{originallemma}{Original lemma}{Original lemmas}
\makeatletter
\newlabel{etale-cohomology-lemma-connected-locally-constant}{{09BF}{source}{}{sourcepointer-1}{}}
\newlabel{etale-cohomology-lemma-connected-locally-constant@cref}{{[originallemma][0][]09BF}{[1][0][]source}{}{}{}}
\makeatother
\crefname{originalreference}{original reference}{original references}
\Crefname{originalreference}{Original reference}{Original references}
\makeatletter
\newlabel{lemma-compare-locally-constant}{{099Y}{source}{}{sourcepointer-2}{}}
\newlabel{lemma-compare-locally-constant@cref}{{[originalreference][0][]099Y}{[1][0][]source}{}{}{}}
\makeatother
\crefname{originalreference}{original reference}{original references}
\Crefname{originalreference}{Original reference}{Original references}
\makeatletter
\newlabel{lemma-weird}{{09BV}{source}{}{sourcepointer-3}{}}
\newlabel{lemma-weird@cref}{{[originalreference][0][]09BV}{[1][0][]source}{}{}{}}
\makeatother
\crefname{originallemma}{original lemma}{original lemmas}
\Crefname{originallemma}{Original lemma}{Original lemmas}
\makeatletter
\newlabel{more-algebra-lemma-isomorphic-completions}{{09BC}{source}{}{sourcepointer-4}{}}
\newlabel{more-algebra-lemma-isomorphic-completions@cref}{{[originallemma][0][]09BC}{[1][0][]source}{}{}{}}
\makeatother
\title{A lisse constructible sheaf on a connected scheme over the pro-etale site is adic lisse}
\author{The Stacks Project Authors}
\date{stacks-project:proetale.tex:local-snapshot-20261012; GNU FDL 1.2 or later}
\begin{document}
\maketitle
\noindent\textit{Editorial note.} All delivered mathematical text is the author's own, copied from the pinned Stacks Project chapter file proetale.tex; no statement or proof step was rewritten, reordered or removed.\par\medskip
\noindent\textit{Editorial note.} The counted proof is the author's complete argument: the identification of the reductions through the comparison with the etale site and connectedness, the local description on a covering by affine weakly contractible schemes, and the comparison of the resulting finite modules through their reductions and completions.\par\medskip
\noindent\textit{Editorial note.} The label lemma-connected-lisse names the counted statement, and the proof environment immediately following it in the source is delivered verbatim as its proof.\par\medskip
\noindent\textit{Editorial note.} The counted statement is carried by the author's own lemma environment in the chapter source.\par\medskip
\noindent\textit{Editorial note.} Pre-open gate: python operations/agent-authoring/tools/precard\_gate.py --source proetale.tex --label lemma-connected-lisse --start 5845 --end 5871 --local-id 3597 returned RESULT=OK (L3597 free, occupancy FREE, public interval clear). Candidate taken from sources/stacks-proetale/candidates-halfA.json.\par\medskip
\noindent\textit{Editorial note.} The statement and proof are copied from the pinned chapter file; the cited inputs are kept as external pointers with their canonical Stacks tags.\par\medskip
\noindent\textit{Allowed standard prerequisites (editorial).}\begin{itemize}
\item Equivalence between locally constant sheaves on the etale and pro-etale sites. Same chapter, the lemma comparing locally constant sheaves of finite presentation on the two sites through the pullback. Supplies the locally constant sheaves G\_n of Lambda/I\textasciicircum{}n-modules whose pullbacks are the reductions F/I\textasciicircum{}n F.
\item Local constancy of a locally constant sheaf on a connected scheme. Etale Cohomology, the lemma that a locally constant sheaf of modules on a connected scheme is locally isomorphic to a constant sheaf on a finite module. Supplies the finite modules M\_n with G\_n locally isomorphic to the constant sheaf on M\_n.
\item Coverings by affine weakly contractible schemes. Same chapter, the lemma producing a covering by affine weakly contractible schemes. Supplies the covering \{W\_t -> X\} on which the restrictions of F can be analysed.
\item Description of lisse constructible sheaves on an affine weakly contractible scheme. Same chapter, the lemma identifying such a sheaf with the completion of a finite module. Supplies the finite completed modules N\_t with F restricted to W\_t isomorphic to the completion of N\_t.
\item Isomorphism of finite modules with isomorphic completions. More on Algebra, the lemma that two finite modules over a Noetherian ring with isomorphic completions are isomorphic. Supplies the conclusion that the N\_t are all isomorphic, which makes F adic lisse.
\end{itemize}\medskip
\ifdefined\hypertarget\hypertarget{sourcepointer-1}{}\fi
\noindent\textit{Source reference.} Original lemma 09BF: sources/stacks-chow/tags.txt, source lines 11793--11793. Prerequisite pointer invoked at proof line(s) 5858; it names a standard fact of a sibling chapter that the counted argument uses, and every step of the counted proof built on it is retained verbatim in the delivered body.\par\medskip
\ifdefined\hypertarget\hypertarget{sourcepointer-2}{}\fi
\noindent\textit{Source reference.} Original reference 099Y: sources/stacks-proetale/pinned-original-source.tex, source lines 4338--4338. Prerequisite pointer invoked at proof line(s) 5854; it names a standard fact of the same chapter that the counted argument uses, and every step of the counted proof built on it is retained verbatim in the delivered body.\par\medskip
\ifdefined\hypertarget\hypertarget{sourcepointer-3}{}\fi
\noindent\textit{Source reference.} Original reference 09BV: sources/stacks-proetale/pinned-original-source.tex, source lines 5783--5783. Prerequisite pointer invoked at proof line(s) 5864; it names a standard fact of the same chapter that the counted argument uses, and every step of the counted proof built on it is retained verbatim in the delivered body.\par\medskip
\ifdefined\hypertarget\hypertarget{sourcepointer-4}{}\fi
\noindent\textit{Source reference.} Original lemma 09BC: sources/stacks-chow/tags.txt, source lines 11790--11790. Prerequisite pointer invoked at proof line(s) 5869; it names a standard fact of a sibling chapter that the counted argument uses, and every step of the counted proof built on it is retained verbatim in the delivered body.\par\medskip
\begin{lemma}
\label{lemma-connected-lisse}
Let $X$ be a connected scheme. Let $\Lambda$ be a Noetherian ring and let
$I \subset \Lambda$ be an ideal. If $\mathcal{F}$ is a lisse
constructible $\Lambda$-sheaf on $X_\proetale$, then $\mathcal{F}$
is adic lisse.
\end{lemma}

\begin{proof}
By Lemma \ref{lemma-compare-locally-constant} we have
$\mathcal{F}/I^n\mathcal{F} = \epsilon^{-1}\mathcal{G}_n$
for some locally constant sheaf $\mathcal{G}_n$ of $\Lambda/I^n$-modules. By
\'Etale Cohomology, Lemma
\ref{etale-cohomology-lemma-connected-locally-constant}
there exists a finite $\Lambda/I^n$-module $M_n$ such that
$\mathcal{G}_n$ is locally isomorphic to $\underline{M_n}$.
Choose a covering $\{W_t \to X\}_{t \in T}$ with each $W_t$
affine and weakly contractible.
Then $\mathcal{F}|_{W_t}$ satisfies the assumptions of
Lemma \ref{lemma-weird}
and hence $\mathcal{F}|_{W_t} \cong \underline{N_t}^\wedge$
for some finite $\Lambda^\wedge$-module $N_t$. Note that
$N_t/I^nN_t \cong M_n$ for all $t$ and $n$. Hence
$N_t \cong N_{t'}$ for all $t, t' \in T$, see
More on Algebra, Lemma \ref{more-algebra-lemma-isomorphic-completions}.
This proves that $\mathcal{F}$ is adic lisse.
\end{proof}

\end{document}
